\documentclass[twocolumn, 12pt]{article}

\usepackage[utf8]{inputenc}
\usepackage[english, spanish]{babel}
\usepackage{fullpage}
\usepackage{graphicx}
\usepackage{amsmath}
\usepackage{enumitem}
\usepackage{chngcntr}
\usepackage{setspace}
\usepackage{url}
\usepackage{csquotes}
\usepackage{float}
\usepackage{verbatim}
\usepackage{tabularx}
\usepackage{amsmath}
\usepackage{caption}
\usepackage{bm}
\usepackage{colortbl}
\usepackage{xcolor}

\usepackage{multirow}

% \usepackage{hyperref}

\definecolor{LigthGray}{rgb}{0.7098, 0.7294, 0.7215}
\definecolor{LigthGrayPlus}{rgb}{0.6862, 0.7254, 0.7294}
\definecolor{White}{rgb}{1, 1, 1}
\definecolor{Red}{rgb}{1, 0, 0}

\counterwithin{figure}{section}
\renewcommand{\thesection}{\arabic{section}}
\renewcommand{\thesubsection}{\thesection.\arabic{subsection}}
\renewcommand{\baselinestretch}{1.5}

\usepackage[style=apa, maxnames=6, minnames=3, backend=biber]{biblatex}
\DefineBibliographyStrings{english}{%chktex-file 1 chktex-file 6
    andothers = {\em et\addabbrvspace al\adddot}
}
\addbibresource{./Bibliography/bibliography.bib}

\usepackage{array}
\usepackage{enumitem}

\setlength{\parskip}{0pt}

\raggedbottom{}

\newcommand{\bolditalic}[1]{\textbf{\textit{#1}}}

\begin{document}

\begin{titlepage}
    \centering
    \includegraphics[width=0.3\textwidth]{Images/logo_utb.png}\par\vspace{1cm}
    {\scshape\LARGE Universidad Tecnológica de Bolívar \par}
    \vspace{1cm}

    {\scshape\Large FÍSICA CALOR Y ONDAS \par}
    \vspace{.2cm}

    % chktex-file 8
    {\scshape\Large Grupo 1 \par}
    \vspace{1cm}
    % chktex-file 8
    \slshape {\Large \bfseries{}Informe de Laboratorio No. I\\}
    \vspace{2cm}

    \slshape {\itshape{} Mauro González, T00067622 \\}
    \slshape {\itshape{} German De Armas Castaño, T00068765 \\}
    \slshape {\itshape{} Angel Vega Rodriguez, T00068186 \\}
    \slshape {\itshape{} Juan Jose Osorio Ariza, T00067316 \\}
    \slshape {\itshape{} Jorge Alberto Rueda Salgado, T00068722 \\}
    \vfill
    Revisado Por \\
    Duban Andres Paternina Verona\\
    {\large \today\par}
\end{titlepage}

% ! ----------------------------------------------------------------------|>
\section{Introducción}

% ! ----------------------------------------------------------------------|>
\section{Objetivos}

% + ----------------------------------------|>
\subsection{Objetivo general}

% + ----------------------------------------|>
\subsection{Objetivos específicos}

% ! ----------------------------------------------------------------------|>
\section{Marco Teórico}

% ! ----------------------------------------------------------------------|>
\section{Montaje Experimental}

% ! ----------------------------------------------------------------------|>
\section{Datos Experimentales}

% + ----------------------------------------|>
\subsection{Péndulo Simple}

% chktex-file 44

\begin{table}[H]
    \begin{tabularx}{\linewidth}{|>{\centering\arraybackslash}X|>{\centering\arraybackslash}X|>{\centering\arraybackslash}X|}
        \hline
        \rowcolor{LigthGray} No.   & Longitud \bolditalic{(M)} & Angulo \bolditalic{(°)} \\ \hline
        1                          & 0.3700                    & 15.0000                 \\\hline
        \rowcolor{LigthGrayPlus} 2 & 0.3050                    & 15.0000                 \\\hline
        3                          & 0.4450                    & 15.0000                 \\\hline
    \end{tabularx}

\end{table}

\vspace{-.5cm}

\begin{table}[H]
    \begin{tabularx}{\linewidth}{|>{\centering\arraybackslash}X|>{\centering\arraybackslash}X|>{\centering\arraybackslash}X|}
        \hline
        \rowcolor{LigthGray} \multicolumn{3}{|c|}{Tiempo \bolditalic{(s)}} \\ \hline
        12.4100                          & 12.5100 & 12.1100               \\ \hline
        \rowcolor{LigthGrayPlus} 11.7000 & 11.1500 & 11.0300               \\\hline
        13.5000                          & 13.2800 & 13.3800               \\\hline
    \end{tabularx}

\end{table}

\vspace{-.5cm}

\begin{table}[H]
    \begin{tabularx}{\linewidth}{|>{\centering\arraybackslash}X|>{\centering\arraybackslash}X|}
        \hline
        \rowcolor{LigthGray} Promedio \bolditalic{(s)} & Periodo \bolditalic{(Hz)} \\ \hline
        12.3433                                        & 1.2343                    \\\hline
        \rowcolor{LigthGrayPlus} 11.2933               & 1.1293                    \\\hline
        13.3867                                        & 1.3387                    \\\hline
    \end{tabularx}

\end{table}

\vspace{-.5cm}

\begin{table}[H]
    \begin{tabularx}{\linewidth}{|>{\centering\arraybackslash}X|>{\centering\arraybackslash}X|}
        \hline
        \rowcolor{LigthGrayPlus} \textbf{Oscilaciones} & 10 \\\hline
    \end{tabularx}
\end{table}

% + ----------------------------------------|>
\subsection{Péndulo Compuesto}

\begin{table}[H]
    \begin{tabularx}{\linewidth}{|>{\centering\arraybackslash}X|>{\centering\arraybackslash}X|>{\centering\arraybackslash}X|>{\centering\arraybackslash}X|}
        \hline
        \rowcolor{LigthGray} No. & Masa \bolditalic{(Kg)} & Longitud \bolditalic{(M)} & Distancia \bolditalic{(M)} \\ \hline
        1                        & 0,0490                 & \multirow{3}{*}{0.2470}   & 0,0500                     \\
        2                        & 0,0490                 &                           & 0,0840                     \\
        3                        & 0,0490                 &                           & 0,1020                     \\\hline
    \end{tabularx}
\end{table}

\vspace{-.5cm}

\begin{table}[H]
    \begin{tabularx}{\linewidth}{|>{\centering\arraybackslash}X|>{\centering\arraybackslash}X|>{\centering\arraybackslash}X|}
        \hline
        \rowcolor{LigthGray} \multicolumn{3}{|c|}{Tiempo \bolditalic{(s)}} \\ \hline
        4,0000                          & 3,7100 & 3,7400                  \\\hline
        \rowcolor{LigthGrayPlus} 3,8400 & 3,7000 & 4,1000                  \\\hline
        3,9000                          & 3,8100 & 3,9800                  \\\hline

    \end{tabularx}
\end{table}

\vspace{-.5cm}

\begin{table}[H]
    \begin{tabularx}{\linewidth}{|>{\centering\arraybackslash}X|>{\centering\arraybackslash}X|}
        \hline
        \rowcolor{LigthGray} Promedio \bolditalic{(s)} & Periodo \bolditalic{(Hz)} \\ \hline
        3,8167                                         & 0,7633                    \\\hline
        \rowcolor{LigthGrayPlus} 3,8800                & 0,7760                    \\\hline
        3,8967                                         & 0,7793                    \\\hline
    \end{tabularx}
\end{table}

\vspace{-.5cm}

\begin{table}[H]
    \begin{tabularx}{\linewidth}{|>{\centering\arraybackslash}X|>{\centering\arraybackslash}X|}
        \hline
        \rowcolor{LigthGrayPlus} \textbf{Oscilaciones} & 5 \\\hline
    \end{tabularx}
\end{table}

% + ----------------------------------------|>
\subsection{Péndulo de Resorte}

\begin{table}[H]
    \begin{tabularx}{\linewidth}{|>{\centering\arraybackslash}X|>{\centering\arraybackslash}X|}
        \hline
        \rowcolor{LigthGray} No.   & Masa \bolditalic{(Kg)} \\\hline
        1                          & 0,0100                 \\\hline
        \rowcolor{LigthGrayPlus} 2 & 0,0150                 \\\hline
        3                          & 0,0200                 \\\hline
    \end{tabularx}
\end{table}

\vspace{-.5cm}

\begin{table}[H]
    \begin{tabularx}{\linewidth}{|>{\centering\arraybackslash}X|>{\centering\arraybackslash}X|}
        \hline
        \rowcolor{LigthGray} Longitud Inicial \bolditalic{(M)} & Longitud Final \bolditalic{(M)} \\\hline
        1                                                      & 0,0100                          \\\hline
        \rowcolor{LigthGrayPlus} 2                             & 0,0150                          \\\hline
        3                                                      & 0,0200                          \\\hline
    \end{tabularx}
\end{table}

\vspace{-.5cm}

\begin{table}[H]
    \begin{tabularx}{\linewidth}{|>{\centering\arraybackslash}X|>{\centering\arraybackslash}X|>{\centering\arraybackslash}X|}
        \hline
        \rowcolor{LigthGray} $\Delta X$ \bolditalic{(M)} & K \bolditalic{(N/m)} & Periodo [Calculado] \bolditalic{(Hz)} \\\hline
        0,0750                                           & 1,3080               & 0,5494                                \\\hline
        0,1100                                           & 1,3377               & 0,6653                                \\\hline
        0,1500                                           & 1,3080               & 0,7769                                \\\hline
    \end{tabularx}
\end{table}

\vspace{-.2cm}

\begin{itemize}[label=$\triangleright$]
    \item K Promedio:~\textcolor{Red}{$1.3179$}
\end{itemize}

\vspace{-.2cm}

\begin{table}[H]
    \begin{tabularx}{\linewidth}{|>{\centering\arraybackslash}X|>{\centering\arraybackslash}X|>{\centering\arraybackslash}X|}
        \hline
        \rowcolor{LigthGray} \multicolumn{3}{|c|}{Tiempo \bolditalic{(s)}} \\ \hline
        5,5400                          & 5,5300 & 5,1700                  \\\hline
        \rowcolor{LigthGrayPlus} 6,1300 & 6,2000 & 6,3600                  \\\hline
        8,1300                          & 7,2600 & 7,9600                  \\\hline
    \end{tabularx}
\end{table}

\vspace{-.5cm}

\begin{table}[H]
    \begin{tabularx}{\linewidth}{|>{\centering\arraybackslash}X|}
        \hline
        \rowcolor{LigthGray} Promedio \bolditalic{(s)} \\ \hline
        5,4133                                         \\\hline
        \rowcolor{LigthGrayPlus} 6,2300                \\\hline
        7,7833                                         \\\hline

    \end{tabularx}
\end{table}

\vspace{-.5cm}

\begin{table}[H]
    \begin{tabularx}{\linewidth}{|>{\centering\arraybackslash}X|>{\centering\arraybackslash}X|}
        \hline
        \rowcolor{LigthGrayPlus} \textbf{Oscilaciones} & 10 \\\hline
    \end{tabularx}
\end{table}

% ! ----------------------------------------------------------------------|>
\section{Análisis de datos}

% ! ----------------------------------------------------------------------|>
\section{Conclusiones}

\end{document}